\section{Printing}\label{sec:printing}

Printing works one column at a time: every printable quantity has its own
flag, and there is no default selection. If you pass no print flag, no table
is printed.
Columns always appear in the same fixed order, no matter what order the flags
were given on the command line.

\begin{verbatim}
CJKQC db.jsonl --filename --el --g --rg
\end{verbatim}

\subsection{How the table is formatted}\label{sec:table-format}

The table is plain text output: a header row of column labels,
then one row per record in the order the run produces. Columns are as wide as
their widest cell, with one space between columns and no trailing whitespace.
Text columns are left-aligned; numeric columns are right-aligned; header
cells align with their columns.

An undefined or absent value prints as a single dash \texttt{-}. A
list-valued column with no values also prints \texttt{-}. Numbers keep fixed
per-column formats, so table columns line up and can be pasted into other
tools.

The fixed column order is:

\begin{center}
\small
\begin{tabular}{@{}llll@{}}
\toprule
\multicolumn{4}{c}{\textbf{Fixed print order}} \\
\midrule
\fld{filename} & \fld{name} & \fld{path} & \fld{status} \\
\fld{program} & \fld{command} & \fld{correction} & \fld{el} \\
\fld{zpe} & \fld{h} & \fld{s} & \fld{g} \\
\fld{elcorr} & \fld{hcorr} & \fld{gcorr} & \fld{pop} \\
\fld{popel} & \fld{popelcorr} & \fld{popgcorr} & \fld{rg} \\
\fld{mass} & \fld{dip} & \fld{dips} & \fld{pol} \\
\fld{pols} & \fld{charge} & \fld{multiplicity} & \fld{rsn} \\
\fld{natoms} & \fld{lf} & \fld{freqs} & \fld{nfreq} \\
\fld{nimag} & \fld{rots} & \fld{components} & \\
\bottomrule
\end{tabular}
\end{center}

\subsection{Identity and source columns}\label{sec:identity-columns}

\begin{longtable}{@{}L{3.1cm}L{2.4cm}L{7.7cm}@{}}
\toprule
\textbf{Flag} & \textbf{Header} & \textbf{What it prints} \\
\midrule
\endhead
\flg{--filename} & \fld{filename} & The record's \fld{name}: the stem of the
source file name. \\
\flg{--name} & \fld{name} & The composition written as alphabetical pairs
without separators, for example \texttt{2am3sa}. Undefined without
components. \\
\flg{--path} & \fld{path} & The record's \fld{source_path}, exactly as
stored. \\
\flg{--status} & \fld{status} & \fld{ok}, \fld{failed}, or \fld{truncated}. \\
\flg{--program} & \fld{program} & The program name and version joined by one
space, for example \texttt{xtb 6.7.1}; the name alone when the version is
absent. \\
\flg{--command} & \fld{command} & The stored command line in double quotes,
because it contains spaces. An absent command prints a bare \texttt{-}. \\
\flg{--correction} & \fld{correction} & Correction source details: the kind, then
the program and version, then each command in double quotes. \\
\flg{--components} & \fld{components} & The resolved composition as pairs with
plus signs, for example \texttt{2am+3sa}. \\
\bottomrule
\end{longtable}

\subsection{Energy columns}\label{sec:energy-columns}

The energy columns share one format: hartree with 10 decimals by default, or
kcal/mol with 4 decimals under \flg{--unit}. The list includes derived
thermochemistry and corrected values; each is explained in
Chapter~\ref{sec:thermo} and Chapter~\ref{sec:corrections}.

\begin{center}
\begin{tabular}{@{}lll@{}}
\toprule
\textbf{Flag} & \textbf{Header} & \textbf{Quantity} \\
\midrule
\flg{--el} & \fld{el} & Electronic energy \\
\flg{--zpe} & \fld{zpe} & Harmonic zero-point energy \\
\flg{--h} & \fld{h} & Enthalpy at \flg{--temp} \\
\flg{--s} & \fld{s} & Entropy term $T\,S$ at \flg{--temp} \\
\flg{--g} & \fld{g} & Gibbs free energy at \flg{--temp} \\
\flg{--elcorr} & \fld{elcorr} & Attached correction energy \\
\flg{--hcorr} & \fld{hcorr} & Enthalpy from the correction energy \\
\flg{--gcorr} & \fld{gcorr} & Gibbs free energy from the correction energy \\
\bottomrule
\end{tabular}
\end{center}

\flg{--unit} is a simple on/off switch: adding it alone converts every
hartree-valued print quantity to kcal/mol, every dipole value to debye,
and every polarizability value to \AA$^3$. It takes no value, and a value
after it is a usage error. Without it, energies print in hartree and
dipoles and polarizabilities in atomic units. Relative
cutoffs (Chapter~\ref{sec:filtering}) already compare in kcal/mol and are not
affected by \flg{--unit}, and the dipole cutoff remains in atomic units.

\subsection{Structure and property columns}\label{sec:property-columns}

\begin{longtable}{@{}L{2.9cm}L{2.5cm}L{7.8cm}@{}}
\toprule
\textbf{Flag} & \textbf{Header} & \textbf{What it prints} \\
\midrule
\endhead
\flg{--rg} & \fld{rg} & Mass-weighted radius of gyration in \AA{}, 3
decimals. \\
\flg{--mass} & \fld{mass} & Total mass in u, 3 decimals. \\
\flg{--dip} & \fld{dip} & Dipole magnitude in atomic units, or debye under
\flg{--unit}, 3 decimals. \\
\flg{--dips} & \fld{dips} & The dipole vector's $x,y,z$ components in atomic
units, or debye under \flg{--unit}, 3 decimals each, comma-separated
without spaces. \\
\flg{--pol} & \fld{pol} & Isotropic polarizability in atomic units, or
\AA$^3$ under \flg{--unit}, 3 decimals. \\
\flg{--pols} & \fld{pols} & The nine polarizability tensor values in atomic
units, or \AA$^3$ under \flg{--unit}, 3 decimals each, comma-separated
without spaces, in the order $xx, xy, xz, yx, yy, yz, zx, zy, zz$. \\
\flg{--charge} & \fld{charge} & Net charge as a whole number. \\
\flg{--multiplicity} & \fld{multiplicity} & Spin multiplicity as a whole
number. \\
\flg{--rsn} & \fld{rsn} & Rotational symmetry number as a whole number. \\
\flg{--natoms} & \fld{natoms} & Geometry atom count. \\
\bottomrule
\end{longtable}

\subsection{Frequency columns}\label{sec:frequency-columns}

\begin{longtable}{@{}L{2.9cm}L{2.5cm}L{7.8cm}@{}}
\toprule
\textbf{Flag} & \textbf{Header} & \textbf{What it prints} \\
\midrule
\endhead
\flg{--lf} & \fld{lf} & The lowest stored vibrational frequency in
cm$^{-1}$, 1 decimal. Raw values, so an imaginary mode prints negative. \\
\flg{--freqs} & \fld{freqs} & Every stored vibrational frequency in stored
order, 1 decimal each, comma-separated without spaces. Raw signs are kept. \\
\flg{--nfreq} & \fld{nfreq} & Number of stored vibrational frequencies. \\
\flg{--nimag} & \fld{nimag} & Number of negative (imaginary) frequencies. \\
\flg{--rots} & \fld{rots} & Rotational constants in cm$^{-1}$, descending,
6 decimals each, comma-separated without spaces. \\
\bottomrule
\end{longtable}

\begin{note}
The frequency columns show the \emph{raw stored values}. They are not
changed by \flg{--makereal}, \flg{--ignoreimag}, or \flg{--fscale}; those
settings affect only the thermochemistry derivation
(Section~\ref{sec:freq-manipulation}). A saddle point therefore still prints
a negative \fld{lf} and \fld{nimag} even when its Gibbs energy is made
defined.
\end{note}

\subsection{Boltzmann population columns}\label{sec:populations}

The \flg{--pop} flag prints the fraction of the thermal ensemble that each
printed structure occupies. It is repeatable and takes an optional level:

\begin{center}
\begin{tabular}{@{}lll@{}}
\toprule
\textbf{Invocation} & \textbf{Column} & \textbf{Energy basis} \\
\midrule
\flg{--pop} or \flg{--pop g} & \fld{pop} & Gibbs free energy at \flg{--temp} \\
\flg{--pop el} & \fld{popel} & Electronic energy \\
\flg{--pop elcorr} & \fld{popelcorr} & Attached correction energy \\
\flg{--pop gcorr} & \fld{popgcorr} & Gibbs free energy from the correction \\
\bottomrule
\end{tabular}
\end{center}

Populations are computed over exactly the records the run prints, grouped by
composition, or into one global group under \flg{--noname}. For a group and a
level with values $v_i$ in hartree,

\begin{equation}
p_i = \frac{\exp\!\left(-\dfrac{v_i - v_{\min}}{k_B T}\right)}
           {\sum_j \exp\!\left(-\dfrac{v_j - v_{\min}}{k_B T}\right)},
\end{equation}

where $v_{\min}$ is the group's minimum defined value and $T$ is
\flg{--temp}. Only status \fld{ok} records with a defined level value
contribute. A record with an undefined value prints \texttt{-}; if no member
of the group defines the level, every member prints \texttt{-}. Populations
are dimensionless probabilities printed with 4 decimals, and \flg{--unit}
never affects them.

\begin{warning}
A space-separated value after \flg{--pop} is consumed by the flag. That means
\texttt{CJKQC} \flg{--pop} \texttt{db.jsonl} tries to use \texttt{db.jsonl} as a population
level and fails. Put \flg{--pop} after the positional inputs, or write
\flg{--pop=g}.
\end{warning}

\subsection{Sections and file output}\label{sec:sections-output}

Several print flags produce a section after the main table instead of a
column:

\begin{itemize}
  \item \flg{--formation} prints the formation-energy section
        (Chapter~\ref{sec:formation});
  \item \flg{--formation-corr} prints the corrected formation section;
  \item \flg{--hydration} and \flg{--hydration-corr} print water
        distributions (Chapter~\ref{sec:solvation});
  \item \flg{--solvation CODE} and \flg{--solvation-corr CODE} print the
        distribution for another solvent.
\end{itemize}

The standard-output order is always: main table, raw formation section,
corrected formation section, hydration/solvation sections (raw before
corrected), then the completion summary. A run with only section flags prints
the section and no main table.

\flg{--movie} writes the final structures into \texttt{movie.xyz}, one xyz
frame per \fld{ok} record that carries a geometry, in the order the table
prints. It is a print flag and a visible outcome on its own. Records without
a geometry, or with a non-\fld{ok} status, are skipped and counted in a
standard-error note. The file is always named \texttt{movie.xyz} in the
current working directory and is subject to the same overwrite confirmation
as \flg{--out}.

\flg{--write-xyz} writes one separate xyz file for each \fld{ok} record that
has a geometry and a usable \fld{filename}. The file is named after the
record: a record whose \fld{filename} is \texttt{1sa-3xstr} becomes
\texttt{1sa-3xstr.xyz} in the current working directory. Each file has the
same one-structure content as the corresponding \flg{--movie} frame. It is a
print flag and a visible outcome on its own, and it can be combined with
\flg{--out} and \flg{--movie}.

Records that are not \fld{ok}, have no geometry, have no name, or have a name
that cannot be used as a file name are skipped and counted in a
standard-error note. Names that would produce the same file, including names
that differ only by upper and lower case, stop the run as a usage error
before anything is written. A record's file also must not have the same name
as \texttt{movie.xyz} in a \flg{--movie} run, or as the \flg{--out} database
in a run that writes one. The existing files are covered by one overwrite
confirmation, and \flg{--fout} forces their replacement.

When \flg{--out} is also used, the database is written first, then the
separate xyz files, then \texttt{movie.xyz}. The completion line is
\texttt{wrote <n> xyz files into the current directory}.

\subsection{A complete example}\label{sec:printing-example}

\begin{verbatim}
CJKQC db.jsonl --filename --program --el --g --lf --nimag --pop
\end{verbatim}

This prints the file stem, program name and version, electronic energy,
Gibbs free energy, lowest frequency, imaginary-mode count, and Gibbs-based
population. The exact widths depend on the data; the layout is illustrated
below with hypothetical values.

\begin{verbatim}
filename  program     el            g             lf     nimag  pop
1dma      xtb 6.7.1   -38.5100000000 -38.4987654321 3421.1 0      0.7321
1dma-2    xtb 6.7.1   -38.5090000000 -38.4970000000 3402.7 0      0.2679
\end{verbatim}

Everything printed here is derived from the database at the run's settings.
Nothing in the table was stored when the database was built.
